\begin{textsample}{POS Dim 2 – global\_south\_2025\_04 – Score 23.00 – global\_south\_2025\_04/123047698.txt}  \label{ex:f2_pos_278}
Israel orders mass evacuation in Rafah amid renewed Gaza \textbf{offensive} Displaced Palestinians, who \textbf{flee} from Rafah, arrive in Khan Younis, Gaza, on Monday, March 31, 2025, after the Israeli military issued sweeping evacuation orders covering most of Rafah AP Israel issued sweeping evacuation orders in Rafah, Gaza ' s \textbf{southernmost} \textbf{city}, on March 31, prompting tens of thousands of Palestinians to \textbf{flee} at a time when Muslims were celebrating Eid ul-Fitr, marking the end of the holy month of \textbf{Ramadan}.
The orders, affecting at least 1.4 lakh people according to Philippe Lazzarini of the UN agency for Palestinian \textbf{refugees} ( UNRWA ), hinted at an impending major ground \textbf{operation}.
Lazzarini described the situation as treating people"like pinballs", with constant military directives upending their lives.
The evacuation covered Rafah and the \textbf{areas} between Rafah and Khan Younis.
The IDF ' s Arabic-language spokesman, Avichay Adraee, asked people to shift out of Al-Shawka, Al-Nasr and the \textbf{neighbourhoods} of Al-Salam, Al-Manara and Qizan an-Najjar.
In @ @ @ @ @ @ @ @ @ @ of the \textbf{area} to be cleared, asking the residents to move to the al-Mawasi \textbf{area} on the \textbf{southern} Strip ' s coast, which is crowded with squalid \textbf{tent} \textbf{camps}.
Families \textbf{fled} north, some on foot with \textbf{children} in hand, others piling belongings onto donkeys or cars.
The latest evacuation orders followed \textbf{intensified} Israeli attacks in Rafah on March 29, with more troops deployed.
In \textbf{northern} Gaza, the IDF claimed to have dismantled a Hamas tunnel and a rocket workshop, killing 50"terrorists".
The last major escalation in Rafah, located close to the Egyptian border, took place last May and left the \textbf{city} in ruins.
Israel seized the Rafah \textbf{crossing}, Gaza ' s only non-Israeli-controlled external link, and a border corridor to curb weapons smuggling.
A ceasefire signed in January this year under American pressure required Israel to leave the corridor.
However, it refused, citing security needs.
Earlier last month, Israel ended a two-month truce with Hamas, resuming air and ground attacks.
It cut off @ @ @ @ @ @ @ @ @ @ two million residents in order to pressure Hamas into accepting new truce terms.
Since March 18, the IDF has targeted Hamas leaders and \textbf{infrastructure}, expanding \textbf{operations} into Rafah and \textbf{northern} Gaza.
Israel has proposed a new ceasefire via mediators in Egypt and Qatar, demanding Hamas release 11 living hostages and 16 bodies in exchange for an unspecified number of Palestinian prisoners.
Fighting would pause for 40 days after the hostages ' release, with further talks planned.
Israel also wants Hamas to disarm and leave Gaza.
These are conditions that were not part of the original ceasefire deal and which Hamas rejects.
Hamas insists on enforcing the January agreement, which promised a lasting ceasefire and total Israeli withdrawal from Gaza in exchange for all remaining hostages ( 59, with 24 believed to be alive ).
Talks, however, stalled after preliminary discussions in February.
The latest evacuation orders have drawn sharp criticism.
The Palestinian presidency called them a violation of international law, warning of dire consequences and holding Israel accountable for rising @ @ @ @ @ @ @ @ @ @ terms, holding Israel responsible for not providing proper conditions such as hygiene, \textbf{health} and safety for the displaced.
Israel said that by ordering the evacuation, it was protecting \textbf{civilians} and accused Hamas of using them as human shields, which goes against international law.
The IDF, meanwhile, admitted to mistakenly firing on \textbf{ambulances} and fire engines in \textbf{southern} Gaza last week, believing those were Hamas vehicles.
More than 1,000 people have been killed in Israel ' s post-truce attacks.
Last week, 14 aid \textbf{workers} ' bodies -- eight from the Palestine Red Crescent Society ( PRCS ), five from civil defence and one from a UN agency -- were found in a Rafah mass grave after Israeli attacks.
One PRCS \textbf{medic} remains missing.
Israel says it aims to limit \textbf{civilian} deaths, but critics, including the UN, argue that its actions fall short of \textbf{humanitarian} standards.

% matched lemmas: ambulance, area, camp, child, city, civilian, crossing, flee, health, humanitarian, infrastructure, intensified, medic, neighbourhood, northern, offensive, operation, ramadan, refugee, southern, southernmost, tent, worker
\end{textsample}
